%%
%% licence       kaneton licence
%%
%% project       kaneton
%%
%% file          /home/mycure/kaneton/view/papers/kaneton/kaneton.tex
%%
%% created       julien quintard   [thu dec  8 00:26:00 2005]
%% updated       julien quintard   [thu feb 23 07:35:20 2006]
%%

%
% class
%

\documentclass[12pt,a4wide]{report}

%
% packages
%

\usepackage[english]{babel}
\usepackage{a4wide}
\usepackage[Lenny]{fncychap}
\usepackage{fancyheadings}
\usepackage{multicol}
\usepackage{indentfirst}
\usepackage{graphicx}
\usepackage{color}
\usepackage{xcolor}
\usepackage{verbatim}
\usepackage{xspace}

\pagestyle{fancy}

\setlength{\footrulewidth}{0.3pt}
\setlength{\parindent}{0.3cm}
\setlength{\parskip}{2ex plus 0.5ex minus 0.2ex}

%
% verbatim font
%

\definecolor{verbatimcolor}{rgb}{0,0.4,0}

\makeatletter
\renewcommand{\verbatim@font}
  {\ttfamily\footnotesize\color{verbatimcolor}\selectfont}
\makeatother

%
% kaneton
%

\newcommand{\kaneton}{kaneton\xspace}

%
% header
%

\rfoot{\scriptsize{The kaneton microkernel project}}
\lhead{\scriptsize{The kaneton microkernel project}}

\date{\scriptsize{\today}}

%
% title
%

\title{\huge{The kaneton micro-kernel reference}
\\\vspace{2cm}\includegraphics{logo.jpg}}

%
% authors
%

\author{\small{Julien Quintard}, \small{Matthieu Bucchianeri},
\small{Renaud Voltz}\vspace{2cm}}

%
% document
%

\begin{document}

%
% title
%

\maketitle

%
% --------- text --------------------------------------------------------------
%

%
% authors
%

\newpage


The \kaneton project was introduced at EPITA by two students:

\begin{itemize}
\item Julien Quintard \footnote{quinta\_j@epita.fr}
\item Jean-Pascal Billaud \footnote{billau\_j@epita.fr}
\end{itemize}


This document especially describes the kaneton reference micro-kernel
developed by:

\begin{itemize}
\item Julien Quintard
\item Matthieu Bucchianeri \footnote{bucchi\_m@epita.fr}
\item Renaud Voltz \footnote{voltz\_r@epita.fr}
\end{itemize}


But many people contributed to this project and we thank them.

%
% toc
%

\tableofcontents

%
% chapters
%

%%
%% licence       kaneton licence
%%
%% project       kaneton
%%
%% file          /home/buckman/kaneton/view/papers/kaneton/1-goals.tex
%%
%% created       matthieu bucchianeri   [mon jan 30 17:09:17 2006]
%% updated       matthieu bucchianeri   [tue feb  7 15:40:50 2006]
%%

%
% goals of the kaneton project
%

\chapter{Goals of the \kaneton project}

\newpage

\paragraph{}
This   document   describes   the  \kaneton   micro-kernel   reference
implementation.

\paragraph{}
The \kaneton micro-kernel is  studied during the EPITA System, Network
and  Security specialization.   The  project comes  with many  courses
leading student to learn about architectures and kernels and to design
and develop important parts of a micro-kernel.

\paragraph{}
From this fact, \kaneton has been designed in a very particular way to
keep the  design as  simpler as possible  to be understandable  by the
students.

\paragraph{}
\kaneton is not based on any previous kernel code: it is written from
nothing. As we  wanted the kernel to be  portable on any architecture,
we paid  particular attention  in dividing the  code into  two precise
sections: the independent code and the architecture-dependent code.

\paragraph{}
Designing this portability  system, we intended to port  the kernel on
many architectures, so we  had to keep the architecture-dependent code
as  small  as possible.   Moreover,  the  only  requirement about  the
hardware is  the minimal amount of  memory, which must be  at least 32
megabytes.

\paragraph{}
As we wanted the project to be very pedagogic, \kaneton source code is
clear  and heavily  commented. Each  manager  or service  has its  own
directory and  headers. Function name  and prototypes are  unified and
explicit. Each function is preceded by a header explaining the task it
performs and detailing the different steps.

\section*{Project history}

\subsection*{Year 2004}

\paragraph{}
The  project  was  introduced  by  students willing  to  learn  kernel
architectures.

\subsection*{Year 2005}

\paragraph{}
The  first complete  kaneton  project was  scheduled  from January  to
December.

\subsection*{Year 2006}

\paragraph{}
The  kaneton project comes  with a  development environment.  Then the
students  have to  write parts  of the  microkernel without  having to
build a development environment from scratch which takes much time.

\paragraph{}
Moreover, courses  were added to the  EPITA specialization curriculum:
microprocessor's architectures, kernels history\ldots{}






% end of include file

%%
%% licence       kaneton licence
%%
%% project       kaneton
%%
%% file          /home/buckman/kaneton/view/papers/kaneton/2-overview.tex
%%
%% created       matthieu bucchianeri   [mon jan 30 17:09:45 2006]
%% updated       matthieu bucchianeri   [mon jan 30 17:52:23 2006]
%%

%
% overview of the kernel
%

\chapter{Overview of the kernel}

\newpage

\paragraph{}
\kaneton micro-kernel is only the \textbf{core} of an operating system. Main task like hardware drivers or user services are implemented as \textbf{servers}. So the micro-kernel only has a few functionalities:\\

\begin{itemize}
\item Memory management
\item Process management
\item Communication
\item Events and I/O
\end{itemize}

\paragraph{}
In  this chapter, we  will describe  briefly these  tasks and  all the
associated modules and managers.

\section{Memory management}

\paragraph{}
Handling  the  memory  --  from  virtual  address  space  to  physical
addressing  --   is  done   by  three  major   managers:  \textbf{as},
\textbf{segment} and \textbf{region}.

\subsection*{as, the address space manager}

\paragraph{}
In  \kaneton, we  call  an  \textbf{address space}  a  list of  memory
location referenced by a task. Each task has its own address space.

\subsection*{segment, the physical memory manager}

\subsection*{region, memory mapping manager}

\section{Process management}

\section{Communication}

\section{Events and I/O}

\section{Other managers}

% end of include file

%%
%% licence       kaneton licence
%%
%% project       kaneton
%%
%% file          /home/buckman/kaneton/view/papers/kaneton/3-devenv.tex
%%
%% created       matthieu bucchianeri   [mon jan 30 17:31:26 2006]
%% updated       matthieu bucchianeri   [tue feb  7 15:46:54 2006]
%%

%
% development environment
%

\chapter{Development environment}

\newpage

\section{Overview}

\section{Makefile interface}

\section{Shell interface}

\section{Machine and user specific profiles}

\section{Prototype tool}

% end of include file

%%
%% licence       kaneton licence
%%
%% project       kaneton
%%
%% file          /home/buckman/kaneton/view/papers/kaneton/4-source.tex
%%
%% created       matthieu bucchianeri   [mon jan 30 17:32:40 2006]
%% updated       matthieu bucchianeri   [fri feb 10 18:27:14 2006]
%%

%
% description of the source tree
%

\chapter{Description of the source tree}

\newpage

This section describes \kaneton source tree in details:

\begin{verbatim}
 check/
 core/
 drivers
 env/
 export/
 libs/
 programs/
 services/
 tools/
\end{verbatim}

\section{check/}

The \textbf{check}  subdirectory containes \kaneton  test suite. Check
chapter 9 -- \textbf{Test suite} for more information.

\section{core/}

This subdirectory is the most important of the project. It contains
all the microkernel code, the bootloader and the bootstrap.

\subsection*{core/bootstrap/}

This directory contains bootstraps for all supported architectures.

\begin{verbatim}
 core/bootstrap/arch/
                     ia32-virtual
                     ultrasparc
                     ...
\end{verbatim}

\subsection*{core/bootloader/}

This directory contains bootloaders for all supported architectures.

\begin{verbatim}
 core/bootloader/arch/
                      ia32-virtual
                      ultrasparc
                      ...
\end{verbatim}

\subsection*{core/kaneton/}

This place containts \kaneton source code. It is divided as shown below:

\begin{verbatim}
 core/kaneton/
              arch/
              as/
              conf/
              debug/
              id/
              region/
              schedule/
              segment/
              set/
              stats/
              task/
              thread/
\end{verbatim}

Notice that there is one subdirectory per manager.

The  \textbf{conf} sudirectory  contains links  to user  profile files
(\textit{conf.c} and \textit{conf.h}).

\textbf{debug} contains  a console driver  (only for the  kernel early
boot, later, a console service  is launched) and a serial driver (used
by the test tools).

All  architecture  dependent code  is  placed  in \textbf{arch}.  This
directory contains subdirectories for each architecture.

\begin{verbatim}
 core/kaneton/arch/
                   ia32-virtual
                   ultrasparc
                   ...
\end{verbatim}

Each   of    these   subdirectories   contains    files   implementing
architecture-dependent code for each manager.

\section{drivers/}

% end of include file

%%
%% licence       kaneton licence
%%
%% project       kaneton
%%
%% file          /home/mycure/kaneton/view/papers/kaneton/5-coding.tex
%%
%% created       matthieu bucchianeri   [mon jan 30 17:32:57 2006]
%% updated       julien quintard   [wed feb  8 01:19:48 2006]
%%

%
% coding style
%

\chapter{Coding style}

\newpage

\paragraph{}
The \kaneton  project developers  try to follow  a coding  style. This
coding style was introduced to normalize the source code, leading to a
more readable source code.

\paragraph{}
Nevertheless, you  can adapt this coding  style to your own  but try to
follow the rules.

\section*{Headers}

\paragraph{}
Each file must start with an header formatted as shown below:

\begin{verbatim}
/*
 * licence       kaneton licence
 *
 * project       kaneton
 *
 * file          /home/buckman/kaneton/core/kaneton/as/as.c
 *
 * created       julien quintard   [fri feb 11 02:23:41 2005]
 * updated       matthieu bucchianeri   [mon jan 30 20:30:57 2006]
 */
\end{verbatim}

\paragraph{}
An emacs  file for generating  and updating this  header automatically
can be found in \textit{tools/emacs/header.el}. Additionally, you need
to set two environment variables:

\begin{itemize}
\item \textbf{EC\_LICENCE} must be set to ``kaneton licence''
\item \textbf{EC\_DEVELOPER} must be set to your first name and last name
\end{itemize}

\paragraph{}
Please, do not use nicknames in headers.

\section*{Naming convenions}

\paragraph{}
To keep the code as clear as possible, there are several convetions on
type, functions and variables naming.

\subsection*{Variable name}

\paragraph{}
Variable  names   must  be  expressed  in  english   with  lower  case
letters. Here are a few rules you are encouraged to follow:

\begin{itemize}
\item \textbf{sz} suffix for variables representing a size
\item \textbf{n} prefix for variables representing a number of objects
\item \textbf{id} suffix for identifiers
\item \ldots
\end{itemize}

\paragraph{}
Remember that the type name is also used to define a variable. No need
to repeat it in the variable name.

\subsection*{Function name}

\paragraph{}
Function names  must be expressed  in english with lower  case letters
and must  be prefixed by  the file name  they are implemented  in. For
example, a function part of the address space manager must be prefixed
by \textbf{as\_}.

\paragraph{}
These  names must be  chosen carefully:  they must  explicitely define
what the function does.

\subsection*{Type name}

\paragraph{}
As variables  and functions, type  names must be expressed  in english
with lower case letters.

\paragraph{}
Here are the prefixes you must use when writing your own types:

\begin{itemize}
\item \textbf{m\_} for managers main structures
\item \textbf{o\_} for  objects (an object is a  structure which first
64-bits are an identifier)
\item \textbf{i\_} for interfaces (an interface is a set of function pointers)
\item \textbf{d\_} for architecture-dependent structures
\item \textbf{s\_} for general purpose structures
\item \textbf{t\_} for basic and general purpose typedefs
\end{itemize}

\paragraph{}
Notice  that \textbf{d\_}  can be  combined with  other  prefixes, for
example \textbf{do\_} for a dependent object.

\section*{Includes}

\paragraph{}
To keep the  code clear and compact, developer only  need to include a
minimal number of header files:

\begin{itemize}
\item \textit{kaneton.h} for the operating system declarations
\item \textit{klibc.h} for the kaneton specific C library
\end{itemize}

\paragraph{}
These files  are in the include  path, so do not  use relative include
path.

\section*{Types}

\paragraph{}
You  may use  as  soon as  possible  standard types:  \textbf{t\_uint8},
\textbf{t\_sint32}, \ldots{}

\section*{Return values}

\paragraph{}
Every function must report whether it successed or failed. So function
return type should be \textbf{t\_error}.

\paragraph{}
A function  will return  \textbf{ERROR\_NONE} on success  and anything
else on error (for example: \textbf{ERROR\_UNKNOWN}).

\section*{Indentation}

\paragraph{}
There are several indentation rules in \kaneton.

\begin{enumerate}

\item  Field names  of structs  and unions  must be  aligned  with the
struct or union name

\begin{verbatim}
struct       s_set
{
  t_setid    setid;
  t_setsz    size;
  t_type     type;
};
\end{verbatim}

or

\begin{verbatim}
typedef struct
{
  o_id       id;
  t_staid    stats;
  t_setid    container;
}            m_as;
\end{verbatim}

\item Defines and variable must be aligned as shown below:

\begin{verbatim}
#define TASK_PRIOR_CORE     230
#define TASK_HPRIOR_CORE    250
#define TASK_LPRIOR_CORE    210

m_task*                     task;
t_tskid                     ktask = ID_UNUSED;
\end{verbatim}

This rule also applies for variables declarations in functions.

\item Function prototype and body should look like this:

\begin{verbatim}
t_error             stats_function(t_staid          staid,
                                   char*            function,
                                   t_stats_func**   f)
{
  t_sint64          slot = -1;
  t_sint64          i;

  ...
}
\end{verbatim}

Notice  that  parameters name  are  aligned  between  each other,  and
variable names are aligned with function name and between each other.

Try  to respect  this alignment  between functions  in a  single file:
function names may be all aligned and parameter names also.

\end{enumerate}

\section*{Comments}

\paragraph{}
As \kaneton  is intended  to be a  pedagogical project with  clear and
understandable source code.  No need to say that  comments take a very
important part of this objective.

\paragraph{}
Every file must  begin with a comment describing what  is done in this
code. And  every function must be  preceded by a  comment defining its
behavior.

\paragraph{}
To prevent direct comments in the code, we used ``steps'':

\begin{itemize}

\item Each critical  code section on a function is  preceded by a step
number

\item The function header comment contains steps descriptions.

\end{itemize}

There is an example below:

\begin{verbatim}
/*
 * this function shows the usage of comments and steps.
 *
 * steps:
 *
 * 1) compute the index.
 * 2) make the operation.
 * 3) check the result.
 */
t_error         test_foobar(int      a,
                            int      b,
                            int*     c)
{
  int           idx;

  /*
   * 1)
   */

  idx = text_make_index(a, b);

  /*
   * 2)
   */

  idx = idx * a + b;

  /*
   * 3)
   */

  if (idx < 0)
    return ERROR_UNKNOWN;

  *c = idx;

  return ERROR_NONE;
}
\end{verbatim}

\section*{File structure}


% end of include file

%%
%% licence       kaneton licence
%%
%% project       kaneton
%%
%% file          /home/buckman/kaneton/view/papers/kaneton/6-indep.tex
%%
%% created       matthieu bucchianeri   [mon jan 30 17:33:29 2006]
%% updated       matthieu bucchianeri   [mon jan 30 17:33:40 2006]
%%

%
% architecture-independant code
%

\chapter{Architecture-independant code}

% end of include file

%%
%% licence       kaneton licence
%%
%% project       kaneton
%%
%% file          /home/buckman/kaneton/view/papers/kaneton/7-tools.tex
%%
%% created       matthieu bucchianeri   [mon jan 30 17:33:49 2006]
%% updated       matthieu bucchianeri   [mon jan 30 17:33:59 2006]
%%

%
% tools
%

\chapter{Tools}

% end of include file

%%
%% licence       kaneton licence
%%
%% project       kaneton
%%
%% file          /home/buckman/kaneton/view/papers/kaneton/8-tests.tex
%%
%% created       matthieu bucchianeri   [mon jan 30 17:34:07 2006]
%% updated       matthieu bucchianeri   [tue feb  7 15:47:24 2006]
%%

%
% test suite
%

\chapter{Test suite}

\newpage

% end of include file

%%
%% licence       kaneton licence
%%
%% project       kaneton
%%
%% file          /home/buckman/kaneton/view/papers/kaneton/9-glossary.tex
%%
%% created       matthieu bucchianeri   [mon jan 30 17:34:32 2006]
%% updated       matthieu bucchianeri   [mon jan 30 17:34:43 2006]
%%

%
% glossary
%

\chapter{Glossary}

% end of include file


%
% epilogue
%

\end{document}
\label{}
\ref{}
\pageref{}
